% !TeX root = ../main.tex
\clearpage
\pagenumbering{arabic}
\ChapterOpening{regular-competition}{1}{常规比赛报销}
\ChapterGuide
  {\GuideItem{\hyperref[invoice-title]{中山大学开票信息}}
   \GuideItem{\hyperref[evidence-chain]{各类凭证的用途}}
   \GuideItem{\hyperref[material-packaging]{票据整理与打包}}}
  {\GuideItem{\hyperref[regular-competition]{常规比赛报销}}}
  {\GuideItem{\hyperref[cross-campus-training]{跨校区训练车费报销}}
   \GuideItem{\hyperref[training-allowance]{训练费发放}}}

\setlength{\parindent}{2\ccwd}

\noindent 对于一般的外出比赛，通常会产生以下几类需要个人垫付的费用：
\begin{itemize}[label=\textcolor{BookPurple}{\textbullet},itemsep=3pt,topsep=4pt]
\item 赛前体检费
\item 比赛期间保险费
\item 从学校到赛场的交通费（可能是往返）
\item 比赛住宿地和赛场之间的通勤费用
\item 比赛期间购买的餐饮、饮用水、运动补给等
\end{itemize}

\begin{SummaryBox}
另外，在实际比赛流程中，还会产生报名费，参会会费，住宿费等，这些大额费用请\strong{务必让教练自行处理}，\strong{不要自己垫付！！}

（如果教练想让你直接付款，请你谢绝，或要求先通过微信转账相当金额后转付）
\end{SummaryBox}

\begin{SummaryBox}
  \SmallHeading{关于比赛中的三餐}
  一般来说，餐饮有两种方式：
\begin{itemize}[leftmargin=2\ccwd,labelsep=.6\ccwd,itemsep=5pt,topsep=4pt]
\item 统一订购比赛方提供的餐饮
\item 自行解决并发放餐补（通常是100/人/天）
\end{itemize}
虽然后者看起来很诱人，但是在实际过程中，貌似\Strong{一次都没有真正收到过。}所以如果有的选的话，尽量还是选择比赛方提供的餐饮。

（当然如果比赛方不提供餐饮，那只能自己解决了）
\end{SummaryBox}

\Subheading{赛事文件提供业务依据}
赛事规程、参赛通知用于说明比赛名称、时间、地点及参赛要求；学校的参赛安排或审批材料用于说明业务关联与授权。此事项在
理论上应该由教练负责完成。一般队员不需要对此进行准备。但如果你很不幸成为了指定负责人，需要你自行准备相关材料，笔者没有参与过相关
流程，无法提供具体指导，请向教练请教。
\\

\noindent 对于一般个人来说，你需要准备的材料如下：

\Needspace{10\baselineskip}
\Subheading{1. 交通}
交通要尽可能闭环，开始出发地和最终回到的最好是同一个校区，如果不是，请附上情况说明。(见下)

\begin{ReminderBox}
\ReminderTitle{注意！！！}
\normalfont\sffamily\fontsize{10}{17}\selectfont
打车尽量\Strong{避免乘坐计程车}，计程车只能现场开票且提交原件（纸质发票不能提交照片作为报销依据），
一旦忘记开票或者没有提交原件就有可能被卡，或者要走票据丢失流程，相当麻烦。所以请务必使用\Strong{能够线上开票的网约车}进行通勤。
\end{ReminderBox}

\Needspace{17\baselineskip}
\Subheading{凭证清单}
\begingroup
\sffamily\fontsize{10}{17}\selectfont
\setlength{\tabcolsep}{8pt}
\renewcommand{\arraystretch}{1.12}
\arrayrulecolor{BookBlue!20}
\begin{tabularx}{\linewidth}{@{}p{24mm}p{57mm}Y@{}}
\toprule
\rowcolor{BookBlue!5}\Strong{费用类别} & \Strong{凭证材料} & \Strong{说明与引用} \\[4pt]\midrule
\Strong{铁路交通} & 铁路电子客票 & 开具方式见序章第\BookRef{invoices}页。 \\[7pt]\midrule
\Strong{飞机} & 机票行程单、发票 & \\[7pt]\midrule
\Strong{打车} & 打车发票、行程单 & 行程单见序章第\BookRef{itinerary}页。 \\[7pt]\midrule
\Strong{体检} & 体检发票、付款记录 & 付款记录见序章第\BookRef{payment}页。 \\[7pt]\midrule
\Strong{保险} & 保险发票、保单 & 保单相关页面的提取方式见序章第\BookRef{insurance}页。 \\[7pt]\midrule
\Strong{其他} & 发票 & 其他材料视具体要求而定。 \\[4pt]\bottomrule
\end{tabularx}\par
\endgroup



\Needspace{230mm}
\Subheading{2. 个人报销汇总表格}
每个人需要手写或打印一张报销汇总表，写明自己的姓名、学号、邮箱、发票张数和各个发票的金额以及总金额，和报销材料一起提交，参考下图。

\DocumentFigure[.7\linewidth][colframe=black!25,colback=white,drop fuzzy shadow=black!25]{reimbursement-summary.pdf}{1.1}{个人报销汇总表示例 \phantomsection\label{summary-example}}

你可以\href{https://github.com/PhSeCl/SYSU-Reimbursement-Handbook/releases/download/v2026.11.05/reimbursement-summary-form.pdf}{\textcolor{BookBlue}{点击这里}}直接下载可填写的表单，也可以自己完成。

当然，你也可以把它作为《行程情况说明》（见下）的封面。


\Needspace{100mm}
\Subheading{3. 行程情况说明}

在实际外出比赛过程中，想让行程严格保持闭环实际上是非常困难的，会出现许多特殊情况，比如说:
\begin{itemize}[label=\textcolor{BookPurple}{\textbullet},itemsep=3pt,topsep=4pt]
  \item 从A校区出发，结束后回到B校区
  \item 从学校出发，比赛结束后直接回家
  \item 没有直达的交通方式，需要第三地中转
  \item 甲乙共同出行，去程由甲打车，回程由乙打车，体现在各自的报销材料中会有不闭环的情况
  \item 赛事方提供比赛和住宿地之间的接驳巴士，但是由于时间不合适等其他原因，选择了自行打车
\end{itemize}
等等诸如此类，为了避免财务$\to$教练$\to$队员$\to$教练$\to$财务的反复询问，建议附上一份个人的行程情况说明，在说明中解释：

\begin{itemize}[label=\textcolor{BookPurple}{\textbullet},itemsep=3pt,topsep=4pt]
  \item 总体动线以及必要的解释
  \item 交通方式和时间
  \item 费用金额
  \item 非闭环行程的原因说明
\end{itemize}
等等，尽量把所有可能被问到的问题都在说明中解释清楚。
可以参考下图的示例

\par\noindent\begin{minipage}{\linewidth}
  \begin{minipage}[t]{.48\linewidth}
    \DocumentFigure{itinerary-explanation-2.png}{1.2}{行程情况说明示例（第 2 页）}
  \end{minipage}\hfill
  \begin{minipage}[t]{.48\linewidth}
    \DocumentFigure{itinerary-explanation-3.png}{1.3}{行程情况说明示例（第 3 页）}
  \end{minipage}
\end{minipage}



\heading{总结}

常规比赛报销需要把个人垫付的费用、对应的票据和实际行程整理清楚。报名费、参会会费、住宿费等大额费用交\Strong{教练处理}；个人垫付的体检、保险、交通等费用，则需要从付款时就开始留存材料。

整理时，可以先按费用类别归档，再对照本章的凭证清单补齐材料。同一笔费用的发票、付款记录或行程单应放在一起，方便核对。交通尽量形成闭环；涉及跨校区、回家、中转、共同出行等情况时，在行程情况说明中交代实际路线、时间、金额及原因。

\begin{SummaryBox}
\SummaryTitle{提交前需要核对的内容}
\begin{itemize}[leftmargin=2\ccwd,labelsep=.6\ccwd,itemsep=6pt,topsep=4pt]
  \item \Strong{开票信息准确：}购买方为中山大学，纳税人识别号等信息与序章所列一致。
  \item \Strong{凭证材料齐全：}各类费用附上对应材料；纸质发票提交原件，不能只交照片。
  \item \Strong{汇总金额一致：}汇总表写明姓名、学号、邮箱、发票张数、各张发票金额及总额，逐项核对，避免重复或遗漏。
  \item \Strong{特殊行程有说明：}对照票据检查行程，补充需要解释的部分，并写清相关费用的承担方式。
  \item \Strong{文件整理完毕：}电子材料统一为 PDF，使用可读的文件名，分类后连同汇总表和必要的情况说明一起打包提交。
\end{itemize}
\end{SummaryBox}

\begin{SideNote}
  \NoteTitle{整理的麻烦$<<$被过问的麻烦}
  格式转化，文件命名，分类打包，虽然麻烦，但是比起被财务、教练反复询问，要求补充材料或者解释要简单得多。\\
  报销时按照一个赛事作为一个整体来报销的，一个人被卡，所有人都无法下账，麻烦占用您宝贵的一段时间自行认真整理检查，对大家都好\textasciitilde{}\textasciitilde{}\\
  如果拖得太久，第二年可能导致超过报销期限而\strong{无法报销}，所以在比赛后请尽快整理好材料，需要补充时请及时补充。
\end{SideNote}
